\section{Bloom Filters}

\epigraphwidth=0.75\textwidth
\epigraph{\emph{The Ministry of Peace concerns itself with war, the Ministry
of Truth with lies, the Ministry of Love with torture and the Ministry of
Plenty with starvation. These contradictions are not accidental, nor do they
result from from ordinary hypocrisy: they are deliberate exercises in
\textbf{doublethink}.}}{---George Orwell, 1984}

\noindent The Internet is very large, very fast, and full of
\emph{badthink}. The masses spend their days sending each other cat
videos and communicating through their beloved social media platforms:
Facebook, Instagram, Snapchat, Twitter, Discord and worst of all,
TikTok. Some of these will ban users for \emph{badthink} that is
contrary to their political views, but that is not enough for you. To
your dismay, you find that a portion of the masses frequently use words
you deem improper: \emph{oldspeak}. You decide, as the newly elected
Dear and Beloved Leader of the Glorious People's Republic of Santa Cruz
(GPRSC), that a more neutral \emph{newspeak} is required to keep your
citizens of the GPRSC content, pure, and from thinking too much. But how
do you process and store so many words as they flow in and out of the
GPRSC at $10\,$Gbits/second? The solution comes to your brilliant and
pure mind---a \emph{Bloom filter}.

\begin{wrapfigure}{l}{0.185\textwidth}
\centering
        \includegraphics[width=0.18\textwidth]{images/umbrudge.png}
\end{wrapfigure}

A Bloom filter is a space-efficient probabilistic data structure,
conceived by Burton H. Bloom in 1970, and is used to test whether an
element is a member of a set. False-positive matches are possible, but
false negatives are not---in other words, a query for set membership
returns either ``possibly in the set'' or ``definitely not in the set.''
Elements can be added to the set but not removed from it; the more
elements added, the higher the probability of false positives.

\begin{wrapfigure}{r}{0.185\textwidth}
\centering
\includegraphics[width=0.18\textwidth]{images/1984-Big-Brother.jpg}
\end{wrapfigure}

A Bloom filter can be represented as an array of $m$ bits, or a
\textbf{bit vector}. A Bloom filter should utilize $k$ different hash
functions. Using these hash functions, a set element added to the Bloom
filter is mapped to at most $k$ of the $m$ bit indices, generating a
uniform pseudo-random distribution. Typically, $k$ is a small constant
which depends on the desired false error rate $\epsilon$, while $m$ is
proportional to $k$ and the number of elements to be added.

Assume you are adding a word $w$ to your Bloom filter and are using
$k=3$ hash functions, $f(x)$, $g(x)$, and $h(x)$. To add $w$ to the
Bloom filter, you simply set the bits at indices $f(w)$, $g(w)$, and
$h(w)$. To check if some word $w'$ has been added to the same Bloom
filter, you check if the bits at indices $f(w')$, $g(w')$, and $h(w')$
are set. If they are all set, then $w'$ has \emph{most likely} been
added to the Bloom filter. If \emph{any one} of those bits was cleared,
then $w'$ has definitely \emph{not} been added to the Bloom filter. The
fact that the Bloom filter can only tell if some word has \emph{most
likely} been added to the Bloom filter means that \emph{false positives}
can occur. The larger the Bloom filter, the lower the chances of getting
false positives.

So what do Bloom filters mean for you as the Dear and Beloved Leader?
It means you can take a list of proscribed words, \emph{oldspeak} and
add each word into your Bloom filter. If any of the words that your
citizens use seem to be added to the Bloom filter, then this is very
ungood and further action must be taken to discern whether or not the
citizen did transgress. You decide to implement a Bloom filter with
\emph{three} salts for \emph{three} different
hash functions. Why? To reduce the chance of a \emph{false positive}.

You can think of a ``salt'' as an initialization vector or a key. Using
different salts with the same hash function results in a different,
unique hash. Since you are equipping your Bloom filter with three
different salts, you are effectively getting three different hash
functions: $f(x)$, $g(x)$, and $h(x)$. Hashing a word $w$, with
extremely high probability, should result in $f(w) \ne g(w) \ne h(w)$.
These salts are to be used for the SPECK cipher, which requires a
128-bit key, so we have used MD5 \footnote{Rivest, R.. “The MD5
Message-Digest Algorithm.” RFC 1321 (1992): 1-21.} ``message-digest'' to
reduce three books down to 128 bits each. These salts are provided for
you in \texttt{salts.h}. \textcolor{red}{Do not change them.} You will
use the SPECK cipher as a hash function, which will be discussed in
\S\ref{speck}.

The following \texttt{struct} defines the \texttt{BloomFilter} ADT. The
three salts will be stored in the \texttt{primary}, \texttt{secondary},
and \texttt{tertiary} fields. Each salt is 128 bits in size. To hold
these 128 bits, we use an array of two \texttt{uint64\_t}s.

\begin{clisting}{}
struct BloomFilter {
  uint64_t primary[2];   // Primary hash function salt.
  uint64_t secondary[2]; // Secondary hash function salt.
  uint64_t tertiary[2];  // Tertiary hash function salt.
  BitVector *filter;
};
\end{clisting}

This \texttt{struct} definition \emph{must} go in \texttt{bf.c}.

\begin{funcdoc}{\texttt{BloomFilter *bf\_create(uint32\_t size)}}
  The constructor for a Bloom filter. The primary, secondary, and
  tertiary salts that should be used are provided in \texttt{salts.h}.
  Note that you will also have to implement the bit vector ADT for your
  Bloom filter, as it will serve as the array of bits necessary for a
  proper Bloom filter. Bit vectors will be discussed in
  \S\ref{bitvector}.
\end{funcdoc}

\begin{funcdoc}{\texttt{void bf\_delete(BloomFilter **bf)}}
  The destructor for a Bloom filter. As with all other destructors, it
  should free any memory allocated by the constructor and null out the
  pointer that was passed in.
\end{funcdoc}

\begin{funcdoc}{\texttt{uint32\_t bf\_size(BloomFilter *bf)}}
  Returns the size of the Bloom filter. In other words, the number of
  bits that the Bloom filter can access. Hint: this is the length of the
  underlying bit vector.
\end{funcdoc}

\begin{funcdoc}{\texttt{void bf\_insert(BloomFilter *bf, char *oldspeak)}}
  Takes \texttt{oldspeak} and inserts it into the Bloom filter. This
  entails hashing \texttt{oldspeak} with each of the three salts for
  three indices, and setting the bits at those indices in the underlying
  bit vector.
\end{funcdoc}

\begin{funcdoc}{\texttt{bool bf\_probe(BloomFilter *bf, char *oldspeak)}}
  Probes the Bloom filter for \texttt{oldspeak}. Like with
  \texttt{bf\_insert()}, \texttt{oldspeak} is hashed with each of the
  three salts for three indices. If all the bits at those indices are
  set, return \texttt{true} to signify that \texttt{oldspeak} was most
  likely added to the Bloom filter. Else, return \texttt{false}.
\end{funcdoc}

\begin{funcdoc}{\texttt{uint32\_t bf\_count(BloomFilter *bf)}}
  Returns the number of set bits in the Bloom filter.
\end{funcdoc}

\begin{funcdoc}{\texttt{void bf\_print(BloomFilter *bf)}}
  A debug function to print out the bits of a Bloom filter. This will
  ideally utilize the debug print function you implement for your bit
  vector.
\end{funcdoc}
